\documentclass{article}
\usepackage{graphicx} % Required for inserting images
\usepackage{datetime} % Add this line to include the datetime package
\usepackage{amsmath}
\usepackage{amssymb}
\usepackage{amsthm}
\usepackage{tikz}
%\usepackage{hyperref}
\usepackage{xcolor} % Required for defining custom colors

\definecolor{darkblue}{rgb}{0.0, 0.0, 0.5} % Define a darker blue color

\usepackage[colorlinks=true, linkcolor=darkblue, urlcolor=darkblue, citecolor=darkblue]{hyperref} % Use the darker blue color

% Define theorem style
\theoremstyle{plain} % Bold title, italic body
\newtheorem{theorem}{Theorem}
\newtheorem{lemma}[theorem]{Lemma}
\newtheorem{corollary}[theorem]{Corollary}

\theoremstyle{definition} % Bold title, normal body
\newtheorem{definition}[theorem]{Definition}
\newtheorem{example}[theorem]{Example}

\newcommand{\bbZ}{\mathbb{Z}}


\title{Homework 6}
\author{Aaron Honculada}
\date{\monthname[\month] \the\year}

\begin{document}

\maketitle

\section*{Problem 1}
\noindent(a) How many abelian group of order $2025000 = 2^3 \cdot
3^4 \cdot 5^5$ are there up to isomorphism?

We know that the number of unique abelian groups is equivalent
to the integer partitions of the prime powers in the decomposition.
Therefore given $p(n)$ where $p$ is the partition function which
counts the number of integer partitions for a given positive integer.
Thus we have $p(3)=3$, $p(4)=5$, and $p(5)=7$. So the number of 
abelian groups of order $2025000 = 3 \cdot 5 \cdot 7 = \boxed{105}$.

\noindent(b) For which positive integers $n$ does there exist
a unique abelian group of order $n$ up to isomorphism?

There exists a unique abelian group of order $n$ up to isomorphism
for all $n \in \mathbb{Z}_{>0}$ such that $n$ is a square-free integer.

\section*{Problem 2}
\noindent Find the elementary divisors and invariant factors of
the group $U_{19019}$. (Note: $19019 = 7\cdot 11 \cdot 13 \cdot 19$)

To find the order of $U_{19019}$, we first find $\phi(19019)$.
Given the prime decomposition we know $\phi(19019) = (7-1)(11-1)(13-1)(19-1)$
by Euler's Theorem. Thus the order is 12960. Furthermore, since
$(7,11,13,19) = 1$, then by the Chinese Remainder Theorem 
$U_{19019} \cong U_{7} \otimes U_{11} \otimes U_{13} \otimes U_{19}$.
We know $U_7 \cong \bbZ_6$ since the order of $U_7=(7-1)=6$ and
there is an element in $U_7$ with order 6, namely $3\in U_7$.
We make a similar argument for $U_{11} \cong \bbZ_{10}$ since $U_{11}$ has order
$(11-1)=10$ and element $7\in U_{11}$ has order 10. We claim that
$U_{13}$ is isomorphic to $\bbZ_{12}$ which must be isomorphic to
$\bbZ_3 \oplus \bbZ_4$ instead of $\bbZ_3 \oplus \bbZ_2 \oplus \bbZ_2$
because the maximum order of an element in a direct product group is
the least common multiple of the orders of the groups. The LCM$(3,2,2)=6$ while
LCM$(3,4)=12$. Therefore there is an element in
$\bbZ_3 \oplus \bbZ_4$ with order 12 and $\bbZ_3 \oplus \bbZ_4 \cong \bbZ_{12} \cong U_{13}$.
We argue similarly for $U_{19} \cong \bbZ_{18} \cong \bbZ_2 \oplus \bbZ_9$ since the LCM$(2,9)=18$ while
LCM$(2,3,3)=6$.

Given that $U_{19019}$ can be expressed as the direct sum of
the following cyclic groups:
\begin{itemize}
    \item $\bbZ_{2} \oplus \bbZ_{3}$
    \item $\bbZ_{2} \oplus \bbZ_{5}$
    \item $\bbZ_{3} \oplus \bbZ_{4}$
    \item $\bbZ_{2} \oplus \bbZ_{9}$
\end{itemize}
we enumerate the elementary divisors of $U_{19019}$:

$\boxed{2,2,2,3,3,4,5,9}$

We find the invariant factors of $U_{19019}$ by creating
the following table where each row is descending powers of
a prime:
\begin{center}
    \begin{tabular}{c|c|c|c}
        2 & 2 & 2 & 4 \\
          & 3 & 3 & 9 \\
          &   &   & 5 \\
    \end{tabular}
\end{center}
and taking the product of each row gives the order of a
cyclic group of an inviariant factor. Thus the invariant
factors of $U_{19019}$ are $\boxed{2, 6, 6, 180}$ and
$U_{19019} \cong \bbZ_2 \oplus \bbZ_6 \oplus \bbZ_6 \oplus \bbZ_{180}$.

\section*{Problem 3}
\noindent\textit{Let $G$ be a finite abelian group and let $m$
be the largest invariant factor of $G$}

\noindent (a) Show that $m$ is the largest order of an element
of $G$, and the order of any element of $G$ divides $m$.

\begin{proof}
    Given that $m$ is the largest invariant factor of $G$, we know that
    for the other invariant factors $m_1, \dots, m_k$ that $m_i \mid m$ for 
    $1 \leq i \leq k$. Hence, LCM$(ord(m_1),\dots, ord(m_k), ord(m)) = m$ and therefore
    the maximum order of an element in $G$ is $m$ and any element
    in $G$ has order dividing $m$ because it is a subgroup of
    an element of the LCM of the orders of the invariant factors. 
\end{proof}

\noindent (b) Let $H$ be a subgroup of $G$, and let $m'$ and $m''$
be the largest invariant factors of $H$ and $G/H$, respectively.
Show that $m'$ and $m''$ both divide $m$, while $m$ divides 
$m'\cdot m''$.

\begin{proof}
    First we show that both $m' \mid m$ and $m'' \mid m$. Since $H$ is a
    subgroup of $G$ and we know the invariant factors of subgroups
    divide the group, the largest invariant factor $m'$ of $H$ divides $m$.
    Similarly, the invariant factors of the quotient $G/H$ also divide
    those of the ambient group, therefore $m'' \mid m$.

    Finally we show $m \mid m' \cdot m''$. Observe that
    LCM$(m', m'') = m$. Thus $m \mid \text{LCM}(m', m'') \leq m' \cdot m''$.
\end{proof}

\section*{Problem 4}
\noindent Let $G$ be a finite abelian group of order $n$. Show
that there are exactly $n$ homomorphisms from $G$ to $\mathbb{C}^*$.

\begin{proof}
   By the Fundamental Theorem of Finite Abelian Groups, we are able to
   write $G$ as the direct sum of cyclic groups, each of prime
   power order. We define the homomorphism $\phi: G \to \mathbb{C}^*$ which
   sends $k \mapsto e^{\frac{2\pi\cdot i\cdot k}{n}}$ which sends
   each $k \in \bbZ_{p^x}$, where $p^x$ is a prime power, to the $k^{th}$ rooth unity. 
   For each $\bbZ_{p^x}$ we have $p^x$ choices to where send the generator 1. The
   product of all $p^x = n$ therefore there are exactly $n$ homomorphisms
   from $G$ to $\mathbb{C}^*$. 
\end{proof}

\section*{Problem 5}
Let $G$ be a finite abelian group, and let $n$ be a divisor of
the order of $G$. Show that the number of elements of $G$ whose
order divides $n$ is a multiple of $n$.

\begin{proof}
   We observe that for any integer $n \mid |G|$ the number of 
   element that have order $n$ is a multiple of $\phi(n)$. 
\end{proof}

\end{document}